% =====================================================================
\phantomsection
\addcontentsline{toc}{part}{More Languages an Engineer Meets}
\BookChapter[More Languages an Engineer Meets]{More Languages an Engineer Meets}
\label{ch:others}

\begin{researchquestion}
Which other languages will you meet in a database, a factory, a circuit, a workshop or a spreadsheet?
\end{researchquestion}
\OrientationNote
% =====================================================================

The languages in the previous chapters are general-purpose or scientific languages. An
engineer also meets many others: newer systems languages that are replacing C and C++,
languages for databases and automation, languages that run machines in factories,
languages that describe electronic hardware, and languages for modelling physical
systems. This chapter introduces them briefly. Each section answers three questions: what
is it, where will you meet it, and what does it look like?

% ---------------------------------------------------------------------
\LearningOutcomes{Recognise code written in Rust, Go, SQL, Bash, VBA, Structured Text, Verilog, VHDL, Modelica and G-code; say what each language is for and who uses it; write a simple database query; and explain why memory safety, declarative queries and equation-based modelling matter.}

\section{Rust: memory safety without a garbage collector}
\label{sec:rust}

\textbf{Rust} (Mozilla Research, 2010; version 1.0 in 2015) is a systems programming
language that aims at the speed and control of C and C++ while making memory errors ---
buffer overflows, use-after-free, data races between threads --- impossible in ordinary
code. It achieves this with an \term{ownership} system: every value has exactly one owner,
references to it are \emph{borrowed} under rules that the compiler checks, and memory is
freed when the owner goes out of scope. Programs that break the rules do not compile.
\codefile{Rust}{rust/drag.rs}{the drag problem in Rust}
\begin{lstlisting}[style=shell]
$ rustc -O drag.rs && ./drag          # or, in a project:  cargo run --release
\end{lstlisting}
\outfile{rust/drag.out}
Variables are immutable unless declared \code{mut}; types are inferred but strict
(\code{n as f64} converts an integer explicitly); and \code{println!} is a macro with
\py-like format strings. Rust's build tool and package manager \textbf{Cargo} is
widely admired.

\emph{Where you meet it:} operating systems (parts of Linux, Windows and Android), web
browsers, embedded firmware, safety-critical and security-sensitive software, and
command-line tools. Some fast \py tools are written in Rust (the package manager
\code{uv}, the linter \code{ruff}, the data-frame library Polars). Government agencies in
several countries now recommend memory-safe languages such as Rust for new systems
software.

% ---------------------------------------------------------------------
\section{Go: simple, fast, concurrent}

\textbf{Go} (Google, 2009; designed by Robert Griesemer, Rob Pike and Ken Thompson ---
the co-creator of Unix) is a deliberately small, statically typed, compiled language with
garbage collection and built-in support for concurrency (\emph{goroutines} and
\emph{channels}). It compiles quickly to a single executable file with no dependencies.
\codefile{Go}{go/drag.go}{the drag problem in Go}
\outfile{go/drag.out}
\code{:=} declares and initialises a variable with an inferred type, \code{make} creates a
\emph{slice} (a growable array), and the C-like \code{printf} verbs are reused.

\emph{Where you meet it:} cloud infrastructure (Docker, Kubernetes), network services,
command-line tools, and back ends of industrial Internet-of-Things platforms.

% ---------------------------------------------------------------------
\section{SQL: the language of databases}

Engineering data increasingly live in databases: laboratory information systems,
maintenance records, production logs, sensor time series. \textbf{SQL} (Structured Query
Language, IBM 1974, ISO standard since 1987) is the language for defining, querying and
updating \emph{relational} databases, in which data are stored in tables of rows and
columns. SQL is \emph{declarative} (Chapter~\ref{ch:intro}): you describe the result you
want and the database works out how to compute it efficiently.

The following script uses \textbf{SQLite}, a small database engine contained in a single
file, which is built into \py (\code{import sqlite3}), every smartphone and every web
browser. It loads the tensile data and asks two questions.
\codefile{SQL}{sql/tensile.sql}{querying the tensile data in SQL}
\begin{lstlisting}[style=shell]
$ sqlite3 -header -column :memory: < tensile.sql
\end{lstlisting}
\outfile{sql/tensile.out}
(The \code{.import} line is a command of the SQLite shell, not standard SQL.) The core
statements are \code{SELECT \ldots\ FROM \ldots\ WHERE \ldots\ GROUP BY \ldots\ ORDER BY},
\code{INSERT}, \code{UPDATE}, \code{DELETE}, \code{CREATE TABLE}, and \code{JOIN}, which
combines tables --- for example, a table of specimens with a table of material batches.
Other database systems (PostgreSQL, MySQL, Microsoft SQL Server, Oracle) use slightly
different dialects of the same language. pandas, R and every other data-analysis tool can
send SQL queries to a database and receive the results as a table.

% ---------------------------------------------------------------------
\section{Shell scripting: Bash and PowerShell}

The \term{shell} is the program that reads the commands you type in a terminal. Shell
scripts automate sequences of commands: compiling code, running a simulation for many
parameter values, moving and renaming files, submitting jobs to a computing cluster. On
Linux and macOS the standard shell is \textbf{Bash} (or the very similar \textbf{zsh});
on Windows it is \textbf{PowerShell}. Every high-performance computing cluster runs Linux,
so engineers who run large simulations need some Bash.
\codefile{bash}{bash/run_cases.sh}{a parameter study in Bash}
\outfile{bash/run_cases.out}
The script loops over three drag coefficients, creates a modified copy of our \py
program with the stream editor \code{sed}, runs it, extracts the terminal velocity with
\code{grep} and \code{awk}, and stores each result in its own file. The Unix philosophy is
visible here: many small programs, each doing one thing, connected by pipes (\code{|}).
In a real study one would pass the parameter as a command-line argument rather than editing
the source, but this pattern of ``generate input files, run, extract results'' is exactly
how engineers drive commercial and open-source solvers. Job schedulers on clusters (such as
Slurm) are also controlled with Bash scripts.

% ---------------------------------------------------------------------
\section{Excel formulas and VBA}
\label{sec:vba}

The spreadsheet is the most widely used programming environment in engineering, even
though many of its users do not think of it as programming. Excel formulas form a
declarative, functional language: \code{=B2*C2}, \code{=SUM(D2:D50)},
\code{=IF(E2>275,"pass","fail")}, \code{=SQRT(\$B\$1*\$B\$3/\$B\$2)}. Recent versions added
\code{LAMBDA} (user-defined functions), \code{LET} (named intermediate values) and dynamic
arrays, and since 2023 Excel can also run \py code in cells.

\textbf{VBA} (Visual Basic for Applications, 1993) is the macro language built into
Microsoft Office. It is old and Windows-centred, but a vast number of engineering
calculation workbooks in industry depend on it. A VBA user-defined function can be called
from a cell like any built-in function:
\begin{lstlisting}[language=VBA]
' Darcy friction factor (Swamee-Jain), usable in a cell as =FRICTION(B2, C2)
Function FRICTION(Re As Double, relRough As Double) As Double
    If Re < 2300 Then
        FRICTION = 64 / Re
    Else
        FRICTION = 0.25 / (Log(relRough / 3.7 + 5.74 / Re ^ 0.9) / Log(10)) ^ 2
    End If
End Function

' A macro that fills a table with the Euler solution of the drag problem
Sub DragTable()
    Dim m As Double, c As Double, g As Double, dt As Double, v As Double
    Dim n As Integer
    m = 80: c = 0.25: g = 9.81: dt = 0.5: v = 0
    Range("A1:B1").Value = Array("t (s)", "v (m/s)")
    For n = 0 To 40
        Cells(n + 2, 1).Value = n * dt
        Cells(n + 2, 2).Value = v
        v = v + dt * (g - c / m * v ^ 2)
    Next n
End Sub
\end{lstlisting}
Note that VBA's \code{Log} is the natural logarithm, so $\log_{10}$ must be written
\code{Log(x)/Log(10)} --- exactly the kind of detail that causes errors when formulas are
translated between languages.

\begin{engnote}[Engineering note: spreadsheet risk]
Studies of real business spreadsheets have repeatedly found that a large fraction contain
errors, and several costly financial and engineering mistakes have been traced to them.
Spreadsheets are hard to test, formulas are hidden in cells, and copying a formula can
silently shift a reference. For calculations that matter, document the inputs and units
clearly, separate inputs from calculations, lock formula cells, and check results with an
independent method --- or move the calculation into a tested script.
\end{engnote}

% ---------------------------------------------------------------------
\section{PLC languages: Ladder Logic and Structured Text}
\label{sec:plc}

Factories, water-treatment plants, packaging lines and building systems are controlled by
\term{programmable logic controllers} (PLCs): rugged industrial computers that read sensors
and switch motors, valves and heaters in a fixed cycle, typically every few milliseconds.
The international standard \textbf{IEC 61131-3} defines five PLC languages:
\begin{itemize}
\item \textbf{Ladder Diagram} (LD), a graphical language that looks like the relay wiring
  diagrams electricians used before computers: two vertical ``rails'' with horizontal
  ``rungs'' of contacts (inputs) and coils (outputs);
\item \textbf{Function Block Diagram} (FBD), a graphical language of connected blocks;
\item \textbf{Structured Text} (ST), a textual language similar to Pascal;
\item \textbf{Sequential Function Chart} (SFC), for step-by-step sequences;
\item Instruction List (IL), an assembly-like language, now deprecated.
\end{itemize}
Here is a tank-filling controller in Structured Text. The pump starts when the level falls
below a low switch and stops at a high switch (hysteresis again, as in the Arduino example
of Chapter~\ref{ch:c}), and an emergency stop overrides everything:
\begin{lstlisting}[language=StructuredText]
PROGRAM TankControl
VAR_INPUT
    LevelLow    : BOOL;     (* TRUE when the level is below the low switch *)
    LevelHigh   : BOOL;     (* TRUE when the level is above the high switch *)
    EmergencyStop : BOOL;
END_VAR
VAR_OUTPUT
    Pump        : BOOL;
    Alarm       : BOOL;
END_VAR
VAR
    RunTimer    : TON;      (* on-delay timer: alarm if pumping too long *)
END_VAR

IF EmergencyStop OR LevelHigh THEN   (* the conditions that stop the pump come first *)
    Pump := FALSE;
ELSIF LevelLow THEN
    Pump := TRUE;
END_IF;

RunTimer(IN := Pump, PT := T#10m);   (* 10 minutes *)
Alarm := RunTimer.Q;                  (* pump ran > 10 min: possible leak *)
END_PROGRAM
\end{lstlisting}
The same logic as a ladder rung reads: \emph{if} (\code{LevelLow} \emph{or} \code{Pump})
\emph{and not} \code{LevelHigh} \emph{and not} \code{EmergencyStop}, \emph{then} energise
\code{Pump} --- the ``seal-in'' circuit known to every electrician. The order of the tests in
the Structured Text matters. Because the stop conditions are tested first, the pump
switches off at a full tank even if a faulty low-level switch still reports \code{TRUE},
exactly as the ladder rung does. A controller should fail in the safe direction. PLC programs are written
in vendor tools such as Siemens TIA Portal, Rockwell Studio 5000 and CODESYS. Chemical,
process, mechanical and electrical engineers who work in plants will meet them.

% ---------------------------------------------------------------------
\section{Hardware description languages: VHDL and Verilog}

A \term{hardware description language} (HDL) does not describe a sequence of instructions
for a processor; it describes a digital \emph{circuit} --- gates, flip-flops, wires ---
whose parts all operate in parallel. HDL code is \emph{synthesised} into a configuration
for an FPGA (field-programmable gate array) or into the layout of a custom chip (ASIC). The
two main HDLs are \textbf{VHDL} (US Department of Defense, 1980s; Ada-like syntax) and
\textbf{Verilog} (1984; C-like syntax), with its extension SystemVerilog.

A 4-bit counter that counts clock pulses and can be reset:
\begin{lstlisting}[language=Verilog]
module counter4 (
    input  wire       clk,      // clock
    input  wire       reset,    // synchronous reset
    output reg  [3:0] count     // 4-bit output: 0 ... 15
);
    always @(posedge clk) begin // on every rising edge of the clock
        if (reset)
            count <= 4'b0000;
        else
            count <= count + 1; // wraps from 15 back to 0
    end
endmodule
\end{lstlisting}
The same counter in VHDL:
\begin{lstlisting}[language=VHDL]
library ieee;
use ieee.std_logic_1164.all;
use ieee.numeric_std.all;

entity counter4 is
    port (clk, reset : in  std_logic;
          count      : out unsigned(3 downto 0));
end entity;

architecture rtl of counter4 is
    signal c : unsigned(3 downto 0) := (others => '0');
begin
    process (clk)
    begin
        if rising_edge(clk) then
            if reset = '1' then c <= (others => '0');
            else c <= c + 1;
            end if;
        end if;
    end process;
    count <= c;
end architecture;
\end{lstlisting}
Electrical and computer engineers use HDLs to design digital signal processing,
communication interfaces, motor controllers and processors. MATLAB/Simulink (HDL Coder)
and high-level synthesis tools can also generate HDL from C++ or block diagrams.

% ---------------------------------------------------------------------
\section{Modelica: equation-based modelling}
\label{sec:modelica}

In \py or MATLAB you must rearrange a physical model into the form
$\dot{\mathbf{y}} = \mathbf{f}(t, \mathbf{y})$ before you can simulate it. \textbf{Modelica}
(1997, Modelica Association) lets you write the equations as they are, and the tool does
the rearranging. It is \emph{acausal}: an equation \code{m*der(v) = m*g - c*v\^{}2} does not
say which variable is computed from which. Components (resistors, masses, pipes, pumps,
heat exchangers) are connected through \emph{ports}, and the tool assembles and solves the
resulting system of differential--algebraic equations.
\begin{lstlisting}[language=Modelica]
model FallingBody "Falling object with quadratic air drag"
  parameter Real m = 80 "mass (kg)";
  parameter Real c = 0.25 "drag coefficient (kg/m)";
  parameter Real g = 9.81 "gravitational acceleration (m/s2)";
  Real v(start = 0, fixed = true) "speed (m/s)";
  Real h(start = 3000, fixed = true) "height (m)";
equation
  m * der(v) = m * g - c * v^2;
  der(h) = -v;
end FallingBody;
\end{lstlisting}
(\code{start} gives the initial value and \code{fixed = true} tells the tool to treat it
as a firm initial condition, not as a first guess.) Modelica is used for multi-domain system simulation in the automotive, energy, building and
process industries. Tools include the free \textbf{OpenModelica} and the commercial Dymola,
Wolfram System Modeler and others; Modelica's large Standard Library contains thousands of
ready-made mechanical, electrical, thermal and fluid components. Julia's ModelingToolkit
and MATLAB's Simscape follow similar ideas.

% ---------------------------------------------------------------------
\section{G-code: the language of machine tools}

CNC (computer numerical control) milling machines, lathes, laser cutters and 3D printers
are driven by \textbf{G-code} (standardised as ISO 6983 and RS-274). Each line is a command
to move the tool or change a machine state. G-code is usually generated by CAM software
from a CAD model, but machinists read and edit it by hand.
\begin{lstlisting}[language=GCode]
(Cut a 40 mm x 20 mm rectangular pocket outline, 2 mm deep)
G21                (units: millimetres)
G90                (absolute coordinates)
M03 S8000          (spindle on, clockwise, 8000 rpm)
G00 X0 Y0 Z5       (rapid move above the start point)
G01 Z-2 F100       (feed down into the material at 100 mm/min)
G01 X40 F400       (cut along the edges)
G01 Y20
G01 X0
G01 Y0
G00 Z5             (retract)
M05                (spindle off)
M30                (end of program)
\end{lstlisting}

% ---------------------------------------------------------------------
\section{Assembly language}

As explained in Chapter~\ref{ch:intro}, assembly language is a readable form of machine
code, specific to one processor family (x86-64 in most PCs, ARM in phones, Apple
computers and most microcontrollers, RISC-V in a growing number of chips). Engineers rarely
write assembly today, but reading it helps in two situations: debugging low-level embedded
code, and understanding what a compiler has done to your loops. The website Compiler
Explorer (\url{https://godbolt.org}) shows the assembly that C, C++, Fortran, Rust and many
other compilers produce for any code you type, and is an excellent way to learn.

% ---------------------------------------------------------------------
\section{Graphical programming: LabVIEW and Simulink}

Some engineering languages are drawn rather than typed. In \textbf{LabVIEW} (National
Instruments, 1986), programs are \emph{block diagrams}: functions are boxes, and data flow
along wires between them. A box runs as soon as all its inputs have arrived
(\term{dataflow programming}), so parallelism is natural. LabVIEW is widely used for data
acquisition, instrument control and automated test systems in laboratories and production.
\textbf{Simulink} (Chapter~\ref{ch:matlab}) uses block diagrams to describe dynamic systems.
Visual programming is also used in \textbf{Grasshopper} (parametric design in architecture
and structural engineering), \textbf{Node-RED} (Internet of Things) and \textbf{Scratch}
(teaching children). Graphical programs are easy to start but can become hard to manage,
compare and version-control as they grow.

% ---------------------------------------------------------------------
\section{A few more names}

\begin{description}
\item[COBOL] (1959) Business data processing. Hundreds of billions of lines still run in
  banks, insurance companies and governments; COBOL programmers remain in demand.
\item[Pascal, Delphi, Ada] Pascal (1970) was the standard teaching language of the 1980s.
  Its object-oriented descendant, Delphi (Object Pascal), is still used for Windows desktop
  and database applications and is in the TIOBE top 20 (Table~\ref{tab:tiobe}).
  Ada (1980), named after Ada Lovelace, was designed for the US Department of Defense and is
  still used in avionics, railways and other safety-critical systems, often in its
  formally verifiable subset SPARK.
\item[Scratch] (MIT, 2007) A visual language for beginners, especially children. You build a
  program by snapping together coloured blocks (``repeat 10'', ``if \ldots\ then'', ``move 10
  steps'') instead of typing text, so syntax errors are impossible. Many of you may have met
  it at school. The ideas are the same as in \py --- variables, loops, conditions, events
  --- which is why Scratch is a good first step, and why it ranks surprisingly high in
  popularity lists. Block-based programming is also used to program educational robots and
  microcontroller boards (for example the micro:bit).
\item[Dart] (Google, 2011) A C-family language with a syntax close to Java and JavaScript.
  Its main use is \emph{Flutter}, a toolkit with which one program becomes an app for
  Android, iOS, the web and the desktop at the same time. If you ever build a phone app
  for field measurements or inspections, Dart with Flutter is one of the usual choices,
  next to Kotlin and Swift. Hello world is \code{void main() \{ print(\sq{}Hello, world!\sq{}); \}}.
\item[SAS] (1970s) A commercial system and language for statistics and data management,
  long established in the pharmaceutical industry, banking and government statistics. In
  those industries it does the work that R and \py do elsewhere, and regulators are used to
  receiving its output.
\item[Lua] A small, fast scripting language easy to embed in other programs; used in
  games, in some engineering tools and in network equipment.
\item[Kotlin, Swift] Modern languages for Android (Kotlin, on the JVM) and Apple platforms
  (Swift), relevant if you build mobile apps for field engineers.
\item[Haskell, OCaml, F\#, Elixir] Functional languages, valued for correctness and used in
  finance, compilers, formal verification and highly reliable services.
\item[Perl, Ruby, PHP] Scripting languages of the early web. PHP still runs a large share of all websites (for example those built with WordPress). Perl was once the standard
  ``glue'' language of scientific computing, a role \py has taken over.
\item[Mojo] A new language (2023) from the creator of LLVM and Swift, designed as a fast
  superset of \py for AI hardware.
\item[CUDA, OpenCL, SYCL] Extensions and dialects of C/C++ for programming GPUs.
\item[\LaTeX] A typesetting language: these notes are written in it. Its macro language is
  Turing-complete, so in principle it is a programming language too.
\end{description}

\begin{exercises}
\item Install Rust (\url{https://rustup.rs}) and run \file{drag.rs}. Then try to change
  \code{v[n + 1]} inside the loop after removing the \code{mut} keyword from the declaration
  of \code{v}. Read the compiler's error message. How does it help you?
\item Use \py's built-in \code{sqlite3} module to create the \code{tensile} table, insert
  the data from \file{tensile.csv}, and run the two queries of this chapter.
\item Write a Bash script that compiles and runs the C, C++ and Fortran versions of the drag
  program and uses \code{diff} to check that their outputs are identical.
\item Implement the drag table in an Excel sheet using only formulas (one row per time
  step). Compare with the VBA macro. Which would you prefer to give to a colleague, and why?
\item Draw the ladder diagram for the tank controller in this chapter.
\item Install OpenModelica and simulate \code{FallingBody}. When does the body reach the
  ground?
\end{exercises}
